\documentclass[11pt]{letter}
\usepackage[margin=1in]{geometry}
\usepackage{times}
\usepackage[hidelinks]{hyperref}

\setlength{\parindent}{0pt}
\setlength{\parskip}{8pt}

\signature{Ali H. Lazem\\on behalf of both authors}
\address{Ali H. Lazem\\
School of Computer Science and Engineering\\
Bangor University\\
Bangor, UK\\}

\begin{document}

\begin{letter}{The Editors\\
\textit{Nature Medicine}}

\opening{Dear Editors,}

We are pleased to submit our manuscript, ``Do We Need the Largest LLMs for Every Clinical Decision?'', for consideration as an Article in \textit{Nature Medicine}.

The most capable open clinical language models are also the largest, and their hardware demands limit where they can be deployed. A 70--72-billion-parameter model at full precision needs more than 140\,GB of GPU memory, which places it beyond the reach of many hospitals and research groups that, for privacy reasons, would prefer to run models on their own infrastructure. Compression can shrink these models substantially, but it is normally judged by average benchmark scores, which say nothing about whether a smaller model can be trusted on any one clinical decision.

We introduce CERTA, a framework that turns this question into a decision-by-decision, statistically controlled choice. For each decision, CERTA measures how confident the compressed model is and asks whether that decision can be certified to agree with the original full-precision model within a risk level the user sets in advance. Decisions that clear the bar are handled by the small, cheap model; the rest are automatically referred (``escalated'') to the full model. In effect, compression stops being an all-or-nothing setting and becomes an auditable policy: the operator chooses how much disagreement they are willing to tolerate, and the framework guarantees it and flags exactly which decisions the compressed model should not be trusted to make alone.

In plain terms, our main findings are as follows. Across eleven general-purpose and clinically adapted models from three families (7 to 72 billion parameters), tested on 2,500 clinical accept/reject decisions from 2,484 de-identified patients with strictly separated calibration and test sets, 8-bit compression was certified as safe on 99.8 to 100\% of decisions for every model. With a stricter, clinically motivated guarantee applied separately to accept and reject decisions, 4-bit compression let the strongest models run at roughly a quarter of their original size, cutting peak memory from about 146\,GB to about 43\,GB. That is small enough to fit on a single 48\,GB accelerator, rather than the two high-end GPUs the full model requires. Where a group of decisions could not meet the target risk, CERTA escalated that group to the full model instead of issuing an unverified substitution. We also confirmed on two independent, human-annotated disease named-entity benchmarks that the deployed compressed models retain external task performance; importantly, CERTA's guarantee concerns agreement with the reference model, not clinical correctness, and we are careful throughout not to conflate the two.

We believe the study fits \textit{Nature Medicine} because it addresses a concrete barrier to translating capable clinical language models into practice: how to make them small enough for on-premises, resource-constrained, and privacy-sensitive settings while retaining explicit, measurable control over the changes compression introduces. More broadly, it proposes a deployment principle in which computational efficiency is traded against a stated risk level rather than assumed from average benchmark performance, which is directly relevant to the governance of model compression and model updates in clinical workflows.

In the interest of reproducibility and responsible release, all code and artifacts are openly available. The CERTA reference implementation and the complete analysis pipeline, which regenerates every number, table and figure in the manuscript from recorded per-decision data without any model inference, are provided in a public repository with pinned dependencies and a tagged commit. The calibrated per-model certificates and the deployable quantized models are released as a curated Hugging Face collection covering all eleven models, with each artifact packaged under the applicable upstream model licenses. The clinical decisions are derived from the public PMC-Patients corpus; the in-house verification pipeline that generated them is the subject of a separate forthcoming study, and the released per-decision records are sufficient to reproduce every result reported here.

The study is a retrospective computational analysis of PMC-Patients, a publicly available dataset of de-identified patient summaries derived from open-access case reports. No patient contact occurred and no clinical decisions were made or altered as part of the study. The manuscript states its limitations explicitly, including that CERTA controls compression-induced disagreement with the reference model rather than clinical correctness. All authors have approved the manuscript, which is original and is not under consideration elsewhere. We declare no competing interests.

Thank you for considering our manuscript.

\closing{Sincerely,}

\end{letter}

\end{document}
